\scanpage{66}
d) \deleted{\unsure{a}} Les $\Sigma_j\hookrightarrow\Sigma_{ji}$ sont bordants, et si
$x\in\Sigma_i$, \deleted{les} \unsure{dans|donc} les $\Sigma_{ji}\ni x$ (\deleted{pour $j<i$}\inserted{indices}
\deleted{\unsure{tels que}}\deleted{$\Sigma$}\inserted{par l'ens}. $I_{<i,x}\overset{\mathrm{def}}{=}\{j\in I\mid j<i,\ \Sigma_{ji}\ni x\}$,
qui est totalement ordonné par b) sont
des sous-espaces-bords transversaux en $x$
--- i.e. les $\Sigma_{ji}$ $(j\in I_{<i})$ sont des
sous-espaces bords de $\Sigma_i$ \unsure{mutuell\({}^t\)}
transversaux.

e) $\Sigma_{\unsure{i|j}j}\xrightarrow{s_{ij}}\Sigma_i$ est transversal p.r. au système \unsure{trans}-
versal des morphismes bordants
$b_{\unsure{ik|ki}}:\Sigma_{\unsure{k}i}\longrightarrow\Sigma_i$

Si on pose
\[
\Sigma=\coprod_{i\in I}\Sigma_i
\]
\[
\Sigma^{<}(1)=\coprod_{\substack{i,j\in I\\i<j}}\Sigma_{ij}
\]
et $\uline{s},\uline{b}:\Sigma^{<}(1)\mathrel{-}\Sigma$ induits par les
$s_{ij},b_{ij}$, alors les conditions a) b) c)
signifient que $(\uline{s},\uline{b}):\Sigma^{<}(1)\longrightarrow\Sigma\unsure{\times|<}\Sigma$ est
inj. et a comme \deleted{\unsure{graphe}}\inserted{image} le graphe d'une
relation d'ordre \unsure{stricte|strict}, avec $\uline{b}:\Sigma^{<}(1)\longrightarrow\Sigma$
\inserted{immersion} propre, \deleted{et} que l'application \inserted{continue évidente} $\varphi:\Sigma^{\unsure{*}}\unsure{\ell|\deleted{\ell}}\longrightarrow\unsure{\deleted{\Sigma}I|I}$
est strict\({}^t\) croissante, et que pour \unsure{tt|$\forall$} $x\in\Sigma$,
l'application induite $\varphi:\Sigma_{<x}\mathrel{-}I_{<i}$ $(i=\varphi(x))$ est
injective, enfin que $\Sigma_{<x}$ est \unsure{tot.|tt.} ordonné.
\margin{65}
